\section{A non-tree price already forces the square}
\label{sec:tree-prices}

The price-space theorem has a stronger consequence than an embedding
into a space of nonpositive curvature.  The space it constructs has
a price on every finite subset.  Combining that assertion with the
preceding finite tangent-cone limits makes a strict four-point
inequality an exact test for planar flatness.

A \emph{tree pseudometric} is a pseudometric whose quotient by zero
distances embeds isometrically in an $\mathbb R$-tree.  A finite
configuration in such a tree has a finite spanning subtree.  Its
four-point distances obey the familiar rule that the two largest
of the three matching sums are equal.
The scale-free ratio of one matching sum to the larger of the
other two is the quadrilateral constant studied by Dragomir and
Nicas \cite[definition on p.~1042; Theorem 4.1 and Corollary 4.1,
p.~1053]{DragomirNicas2021}. Their CAT(0) upper bound is $\sqrt2$.
Here the price-space realization and finite-limit closure supply
the connection from a strict tree violation to a square price.

\begin{theorem}[A strict non-tree witness is sufficient]
\label{thm:tree-price-dichotomy}
Fix a field $\kk$ and assume
\cite[Proposition C.3, p.~28]{Anthropic2026}.
The following are equivalent:
\begin{enumerate}
\item the unit-square metric is a price over $\kk$;
\item some finite price over $\kk$ is not a tree pseudometric;
\item a four-party price $d$, after relabeling its parties, satisfies
\begin{equation}\label{eq:strict-tree-price-witness}
d_{13}+d_{24}>
\max\{d_{12}+d_{34},d_{14}+d_{23}\};
\end{equation}
\item an ordinary normalized spectral point $\psi$ on four parties satisfies
\begin{equation}\label{eq:strict-tree-spectral-witness}
\psi(\langle2\rangle_{13}\langle2\rangle_{24})>
\max\{\psi(\langle2\rangle_{12}\langle2\rangle_{34}),
       \psi(\langle2\rangle_{14}\langle2\rangle_{23})\}.
\end{equation}
\end{enumerate}
If these conditions fail, every price on every finite party set is a
tree pseudometric.  Over $\C$, that negative branch makes the entire
finite price family equal to the tree pseudometric family.
\end{theorem}

\begin{proof}
For a unit square the three matching sums are $2\sqrt2,2,2$.
It is therefore enough to justify the tree rule, the implication
from a non-tree price to a square price, and the spectral formulation.

For four points in a real tree, take their finite spanning subtree.
After allowing pendant edges of length zero, its shape consists of
four pendant edges and one central connecting segment; the central
segment may also have length zero.  If the pendant lengths sum to
$p$ and the central length is $b$, the three matching sums are
$p,p+2b,p+2b$ in some order.  This follows by counting how often each
edge lies on the two paths in a matching.  The two largest sums are
equal, including when points coincide.  Thus a strict inequality
as in~\eqref{eq:strict-tree-price-witness} is impossible in a tree.

Now let $d$ be a finite price that is not a tree pseudometric.
Proposition C.3 realizes it as a pullback of the metric of a
\emph{price space}.  The unnumbered definition in
\cite[Section C.4, p.~28]{Anthropic2026} requires such a space to be
complete and $\cat$, with its metric a price on every finite subset.
This last requirement is essential here.  The realizing space $X$
cannot be an $\mathbb R$-tree, since $d$ would then be a tree
pseudometric.  Corollary~\ref{cor:closed-metric-planar-family}
applies to the family of prices, which is closed under restrictions,
positive rescaling, and finite pointwise limits.  It shows that
every finite Euclidean planar metric, and in particular the unit
square, is a price.

Finally, Corollary~\ref{cor:price-attained} writes a nonzero finite
price as $d=c\,d_\psi$ for an ordinary normalized spectral point
and $c>0$.  Divide~\eqref{eq:strict-tree-price-witness} by $c$,
exponentiate to base two, and use multiplicativity on the two wires.
This gives~\eqref{eq:strict-tree-spectral-witness}.  Conversely the
wire metric of such a point is itself a price and satisfies the
strict four-point inequality.  The square supplies the witness
for the other direction.

If no square price exists, the implication just proved excludes
every non-tree finite price.  Over $\C$, the converse inclusion of
all tree pseudometrics in the price family is
Theorem~\ref{thm:laminar-tree}.  This last assertion uses that complex
spectral construction; no realization of all trees over an arbitrary
field is needed for the preceding equivalences.
\end{proof}

\begin{corollary}[The diagonal objective has only two values]
\label{cor:tree-price-objective}
Under the same price-space hypothesis, define
\[
R_\square(\kk)=
\max_{\substack{d\in\Price_\kk(\{1,2,3,4\})\\
d_{12},d_{23},d_{34},d_{41}\le1}}(d_{13}+d_{24}).
\]
Then
\[
R_\square(\kk)=
\begin{cases}
2,&\text{if the square is not a price},\\
2\sqrt2,&\text{if the square is a price}.
\end{cases}
\]
\end{corollary}

\begin{proof}
The maximum exists: the defining set is closed, and the triangle
inequality bounds both diagonals by two.  The quadrilateral
inequality~\eqref{eq:price-quad-square} and Cauchy's inequality give
$d_{13}+d_{24}\le2\sqrt2$.  The square realizes equality when it is
a price.

Otherwise Theorem~\ref{thm:tree-price-dichotomy} makes every
admissible $d$ a tree pseudometric.  Its matching sums satisfy
\[
d_{13}+d_{24}\le
\max\{d_{12}+d_{34},d_{14}+d_{23}\}\le2.
\]
The cut metric of $\{1,2\}\mid\{3,4\}$ is a flattening price over
every field.  Its four side values are $0,1,0,1$, and both diagonal
values are one.  It realizes the lower endpoint.
\end{proof}

The theorem changes the positive construction problem.  One need
not prescribe all six distances of a square: any strict violation
of the tree four-point rule by a price suffices.  The strict
inequality is stable under a small perturbation of its six entries.
Neither such a witness nor a proof that every price is a tree
pseudometric is established here.

